\begin{figure*}[t]
\centering
\includegraphics[width=.85\textwidth]{fig06_reward_reliability}
\caption{Bounded reward-reliability evidence.  Panel (a) shows a historical
proxy rising while failure-penalized Vivado utility falls.  Panel (b) shows the
repaired RF midpoint-to-endpoint trajectory: interpolation changes from
\studyclaim{StudyTwoInterpMidFmax} to
\studyclaim{StudyTwoInterpEndFmax}~MHz and extrapolation from
\studyclaim{StudyTwoExtrapMidFmax} to
\studyclaim{StudyTwoExtrapEndFmax}~MHz.}
\label{fig:reward-reliability}
\end{figure*}

\section{Reward diagnostics and validation history}
\subsection{Historical and repaired trajectories}
\label{sec:trajectory}

Fig.~\ref{fig:reward-reliability} contrasts the older exploited reward with
the repaired RF trajectory.  In the historical run, early penalized physical
performance rises from \claim{TrajectoryStartPenalized} to
\claim{TrajectoryEarlyPenalized}~MHz.  Late in training, however, the proxy
rises from \claim{TrajectoryLateProxy} to
\claim{TrajectoryFinalProxy}~MHz while correctness falls from
\claim{TrajectoryLateCorrectness}\% to
\claim{TrajectoryFinalCorrectness}\%.  Penalized physical performance
therefore falls from \claim{TrajectoryLatePenalized} to
\claim{TrajectoryFinalPenalized}~MHz.



A lexical audit identifies one plausible channel in the older reward: a
line-anchored assignment feature increases while its layout-canonical
counterpart does not, and diagnostic rescoring removes
\claim{ReflowRemoved}\% of the raw late reward jump.  The canonical scores are
clamp-saturated and there are \claim{ReflowTokenPairCount} token-identical
cross-checkpoint pairs, so this is bounded attribution rather than a randomized
causal test.  In the repaired run, interpolation remains essentially stable
from midpoint to endpoint and extrapolation improves.  The sealed efficacy study provides the evidence for the repaired reward. The
historical trajectory explains the need for independent physical evaluation.


\subsection{Archived predictor rescoring}
% AUTO-GENERATED by paper/make_revision_evidence.py.
\begin{table}[t]
\centering
\caption{Archived predictor rescoring on the historical optimized distribution. Bias is in MHz; clamp is the number of saturated design cells; rho is cross-design Spearman correlation. Original fits were not rerun for this revision.}
\label{tab:predictors}
\small
\begin{tabular}{@{}lrrr@{}}
\toprule
Configuration & Bias & Clamp & $\rho$ \\
\midrule
\revisionclaim{ArchivePredictorAName} & \revisionclaim{ArchivePredictorABias} & \revisionclaim{ArchivePredictorAClamped} & \revisionclaim{ArchivePredictorARho} \\
\revisionclaim{ArchivePredictorBName} & \revisionclaim{ArchivePredictorBBias} & \revisionclaim{ArchivePredictorBClamped} & \revisionclaim{ArchivePredictorBRho} \\
\revisionclaim{ArchivePredictorCName} & \revisionclaim{ArchivePredictorCBias} & \revisionclaim{ArchivePredictorCClamped} & \revisionclaim{ArchivePredictorCRho} \\
\revisionclaim{ArchivePredictorDName} & \revisionclaim{ArchivePredictorDBias} & \revisionclaim{ArchivePredictorDClamped} & \revisionclaim{ArchivePredictorDRho} \\
\revisionclaim{ArchivePredictorEName} & \revisionclaim{ArchivePredictorEBias} & \revisionclaim{ArchivePredictorEClamped} & \revisionclaim{ArchivePredictorERho} \\
\revisionclaim{ArchivePredictorFName} & \revisionclaim{ArchivePredictorFBias} & \revisionclaim{ArchivePredictorFClamped} & \revisionclaim{ArchivePredictorFRho} \\
\revisionclaim{ArchivePredictorGName} & \revisionclaim{ArchivePredictorGBias} & \revisionclaim{ArchivePredictorGClamped} & \revisionclaim{ArchivePredictorGRho} \\
\bottomrule
\end{tabular}

\end{table}

The archive compares predictors on outputs produced by historical MLP-guided
optimization. Table~\ref{tab:predictors} retains the recorded fits and metrics;
the models were not refitted in the original software environment for this
revision. Several alternatives have smaller bias than the deployed MLP on
this observed distribution, and the character n-gram ridge configuration has
the highest cross-design rank correlation among these archived configurations.
Because this comparison was made after the optimized outputs had been
observed, it measures how well the predictors rescore a fixed distribution.
It does not show how policies trained against those alternatives would behave,
isolate canonicalization in the sealed experiment, or establish resistance to
future reward exploitation. The archive instead documents plausible predictor
alternatives and motivates their prospective evaluation.

\subsection{Executed reward path and contract history}
The RF reward distinguishes
\studyclaim{StudyTwoResolvedGroups} of
\studyclaim{StudyTwoEligibleGroups} eligible RF groups, or
\studyclaim{StudyTwoRewardResolution}\%.  Before sealed generation, all
\studyclaim{StudyTwoContractCandidates} distinct reward-eligible training
candidates reproduced their logged canonical hash and features, compiled, and
passed sampled original/canonical trace comparison; there were
\studyclaim{StudyTwoContractRejected} rejections and
\studyclaim{StudyTwoContractCollisions} cross-design collisions.  The trace
check uses \studyclaim{StudyTwoContractTraceVectors} vectors per frozen oracle
seed and is evidence of the executed reward path, not formal equivalence.
An earlier validation attempt used a stricter contract over every nonempty RF
training completion.  It stopped before sealed efficacy when that contract
encountered a malformed, oracle-rejected completion.  After diagnosing the
scope mismatch, the present study froze a reward-eligible contract over exactly
the candidates that could influence the RF reward, reused the already-trained
checkpoints without retraining or selection, and opened a still-unseen sealed
efficacy split only after the revised gate passed.  The efficacy test is
prospective with respect to that split, but it is not an independent validation
of the post-diagnosis gate scope.

The revision's claim generator recomputes the historical re-anchored outcome
from manifests, oracle summaries, and implementation rows; it revalidates
structural and timing-coverage artifacts before rendering their claims. The
predictor table is an explicitly archival diagnostic. Across both documents,
generated ledgers retain source line ranges and hashes, while the manuscript
verifier checks every included numerical claim and rejects stale outputs.
